\documentclass[11pt,a4paper]{article}

% ---- Packages ----
\usepackage[utf8]{inputenc}
\usepackage[T1]{fontenc}
\usepackage{amsmath,amssymb,amsthm}
\usepackage{mathtools}
\usepackage[margin=1in,headheight=14pt]{geometry}
\usepackage{hyperref}
\usepackage{enumitem}
\usepackage{xcolor}
\usepackage{tcolorbox}
\usepackage{fancyhdr}
\usepackage{booktabs}

% ---- Hyperref ----
\hypersetup{
    colorlinks=true,
    linkcolor=red,
    citecolor=blue,
    urlcolor=blue
}

% ---- Header ----
\pagestyle{fancy}
\fancyhf{}
\fancyhead[L]{\textit{Lesson 12: Repeated Games and Incomplete Information}}
\fancyhead[R]{\thepage}
\renewcommand{\headrulewidth}{0.4pt}

% ---- Theorem Environments ----
\theoremstyle{definition}
\newtheorem{definition}{Definition}[section]
\newtheorem{example}{Example}[section]
\newtheorem{exercise}{Exercise}[section]

\theoremstyle{plain}
\newtheorem{theorem}{Theorem}[section]
\newtheorem{lemma}[theorem]{Lemma}
\newtheorem{proposition}[theorem]{Proposition}
\newtheorem{corollary}[theorem]{Corollary}

\theoremstyle{remark}
\newtheorem{remark}{Remark}[section]

% ---- Colored Boxes ----
\tcbuselibrary{theorems,skins,breakable}

\newtcolorbox{keyresult}[1][]{
    colback=green!5!white,
    colframe=green!50!black,
    fonttitle=\bfseries,
    title=#1,
    breakable
}

\newtcolorbox{rlconnection}[1][]{
    colback=blue!5!white,
    colframe=blue!50!black,
    fonttitle=\bfseries,
    title=#1,
    breakable
}

\newtcolorbox{warning}[1][]{
    colback=red!5!white,
    colframe=red!50!black,
    fonttitle=\bfseries,
    title=#1,
    breakable
}

% ---- Operators ----
\newcommand{\R}{\mathbb{R}}
\newcommand{\N}{\mathbb{N}}
\newcommand{\Q}{\mathbb{Q}}
\newcommand{\Z}{\mathbb{Z}}
\newcommand{\C}{\mathbb{C}}
\newcommand{\Prob}{\mathbb{P}}
\newcommand{\E}{\mathbb{E}}
\newcommand{\Var}{\mathrm{Var}}
\newcommand{\indicator}{\mathbf{1}}
\DeclareMathOperator{\supp}{supp}
\DeclareMathOperator{\conv}{conv}
\DeclareMathOperator{\argmax}{argmax}
\DeclareMathOperator{\argmin}{argmin}
\DeclareMathOperator{\BR}{BR}

% ---- Title ----
\title{\textbf{Lesson 12: Repeated Games and Incomplete Information}\\[10pt]
\large Folk Theorems, Bayesian Games, and Potential Games\\[5pt]
\normalsize Session 12 of 102 $\mid$ Phase I: Mathematical Foundations $\mid$ 8\,h\\
\normalsize Primary text: Fudenberg \& Tirole, \emph{Game Theory}, Chapters 5--6, 8, 12}
\author{Curriculum: A Rigorous Learning Path to Multi-Agent Deep RL}
\date{}

\begin{document}

\maketitle

\begin{abstract}
This document develops the theory of repeated games, games of incomplete information,
and potential games, with complete proofs of all major results. We establish the folk
theorems---characterizing the equilibrium payoff sets of infinitely repeated games---prove
convergence of fictitious play and no-regret dynamics in zero-sum games, develop the
Bayesian game framework and perfect Bayesian equilibrium for signaling games, and
prove that potential games admit pure Nash equilibria with convergent best-response
dynamics. Throughout, we connect these classical results to multi-agent reinforcement
learning: independent learners in repeated games, opponent modeling via Bayesian reasoning,
self-play algorithms, and convergence guarantees in potential games.

\medskip\noindent
\textbf{Prerequisites:} Normal-form games, Nash equilibrium, mixed strategies, minimax
theorem (Lesson 11).

\medskip\noindent
\textbf{Notation:} We write $[n] = \{1, \ldots, n\}$ for player sets. For a strategy
profile $a = (a_1, \ldots, a_n)$, we write $a_{-i} = (a_1, \ldots, a_{i-1}, a_{i+1},
\ldots, a_n)$ and $(a_i', a_{-i})$ for the profile where player $i$ deviates to $a_i'$.
Stage-game payoffs are $g_i(a)$; repeated-game payoffs are $u_i$.
\end{abstract}

\tableofcontents
\newpage

%% ============================================================
%% SECTION 1: Introduction
%% ============================================================
\section{Introduction}
\label{sec:introduction}

In Lesson 11 we studied one-shot (normal-form) games: simultaneous moves, Nash
equilibrium, and the minimax theorem. However, most real strategic interactions---and
virtually all multi-agent reinforcement learning (MARL) scenarios---involve \emph{repeated
interaction}. Agents meet the same opponents across many episodes, observe each other's
past behavior, and adapt their strategies over time. Three fundamental phenomena emerge:

\begin{enumerate}[label=(\roman*)]
    \item \textbf{Cooperation through repetition.} Even when one-shot incentives favor
    defection (as in the Prisoner's Dilemma), the threat of future punishment can sustain
    cooperative outcomes. The \emph{folk theorems} characterize exactly which payoff
    vectors are achievable.

    \item \textbf{Incomplete information.} In practice, agents rarely know each other's
    payoff functions. A robot negotiating with a human does not know the human's utility;
    competing firms do not know rivals' costs. \emph{Bayesian games} model this uncertainty,
    and \emph{perfect Bayesian equilibrium} handles sequential decision-making under such
    uncertainty.

    \item \textbf{Structured interactions.} Many multi-agent problems possess special
    structure---such as \emph{potential games}, where a single global function captures all
    players' incentives. In these settings, simple learning rules (best-response dynamics,
    log-linear learning) are guaranteed to converge.
\end{enumerate}

This document covers all three themes with complete proofs. The roadmap is as follows.
\begin{itemize}
    \item \textbf{Section~\ref{sec:finitely-repeated}}: Finitely repeated games and the
    unraveling argument.
    \item \textbf{Section~\ref{sec:infinitely-repeated}}: Infinitely repeated games, trigger
    strategies, and feasible/individually rational payoffs.
    \item \textbf{Section~\ref{sec:folk-theorems}}: The Nash folk theorem and the subgame
    perfect folk theorem.
    \item \textbf{Section~\ref{sec:zero-sum}}: Fictitious play, no-regret learning, and
    convergence in zero-sum games.
    \item \textbf{Section~\ref{sec:bayesian-games}}: Bayesian games and Bayesian Nash
    equilibrium.
    \item \textbf{Section~\ref{sec:pbe}}: Perfect Bayesian equilibrium and signaling games.
    \item \textbf{Section~\ref{sec:potential-games}}: Potential games, congestion games, and
    convergence of best-response dynamics.
    \item \textbf{Section~\ref{sec:marl-connections}}: Connections to multi-agent RL.
    \item \textbf{Section~\ref{sec:exercises}}: Exercises.
    \item \textbf{Section~\ref{sec:summary}}: Summary of key results.
\end{itemize}

\begin{rlconnection}[Why This Matters for MARL]
In multi-agent RL, agents interact repeatedly and update their policies based on
observed outcomes. The folk theorems explain why independent Q-learners can converge to
cooperative (or adversarial) outcomes that do not arise in single-episode analysis.
Bayesian games formalize the partial observability inherent in opponent modeling.
Potential games identify problem classes where independent learning is guaranteed to
converge---a rare and valuable property in general-sum MARL.
\end{rlconnection}


%% ============================================================
%% SECTION 2: Finitely Repeated Games
%% ============================================================
\section{Finitely Repeated Games}
\label{sec:finitely-repeated}

\subsection{Model}

We begin with a \emph{stage game} $G = (N, (A_i)_{i \in N}, (g_i)_{i \in N})$ where
$N = [n]$ is the player set, $A_i$ is the finite action set of player $i$, and
$g_i : A \to \R$ is the stage-game payoff function, with $A = \prod_{i \in N} A_i$.

\begin{definition}[Finitely Repeated Game]
\label{def:finitely-repeated}
The \emph{$T$-fold repetition} of $G$, denoted $G^T$, is defined as follows:
\begin{enumerate}[label=(\alph*)]
    \item In each period $t \in \{1, 2, \ldots, T\}$, all players simultaneously choose
    actions $a_i^t \in A_i$.
    \item At the end of period $t$, all players observe the realized action profile
    $a^t = (a_1^t, \ldots, a_n^t)$.
    \item The \emph{history} at period $t$ is $h^t = (a^1, a^2, \ldots, a^{t-1}) \in A^{t-1}$.
    The initial history is $h^1 = \emptyset$.
    \item A \emph{(pure) strategy} for player $i$ is a sequence of functions
    $s_i = (s_i^1, s_i^2, \ldots, s_i^T)$ where $s_i^t : A^{t-1} \to A_i$ maps
    histories to actions.
    \item The payoff to player $i$ under strategy profile $s = (s_1, \ldots, s_n)$ is
    \[
        u_i(s) = \frac{1}{T} \sum_{t=1}^T g_i\bigl(a^t(s, h^t)\bigr),
    \]
    where $a^t(s, h^t) = \bigl(s_1^t(h^t), \ldots, s_n^t(h^t)\bigr)$ and histories are
    generated by the strategy profile.
\end{enumerate}
\end{definition}

\begin{remark}
We use average payoffs $\frac{1}{T}\sum_t g_i(a^t)$ rather than total payoffs to keep
the repeated-game payoffs on the same scale as stage-game payoffs, facilitating comparison
with the infinite-horizon discounted case.
\end{remark}

\subsection{Subgame Perfect Equilibrium}

The repeated game $G^T$ is an extensive-form game of perfect information (after each period,
all actions are observed). Thus the natural solution concept is \emph{subgame perfect
equilibrium} (SPE): a strategy profile that constitutes a Nash equilibrium in every subgame.

\begin{definition}[Subgame Perfect Equilibrium]
\label{def:spe}
A strategy profile $s^*$ in $G^T$ is a \emph{subgame perfect equilibrium} (SPE) if, for
every period $t \in \{1, \ldots, T\}$ and every history $h^t \in A^{t-1}$, the
continuation strategy profile $(s_i^{*,\tau}(h^t, \cdot))_{\tau \geq t}$ is a Nash
equilibrium of the subgame starting from $h^t$.
\end{definition}

\subsection{The Unraveling Argument}

The fundamental result for finitely repeated games with a unique stage-game equilibrium is
that repetition adds nothing: the unique SPE plays the stage-game NE in every period.

\begin{theorem}[Unique NE Implies Unique SPE]
\label{thm:unraveling}
If the stage game $G$ has a unique Nash equilibrium $a^* = (a_1^*, \ldots, a_n^*)$, then
the $T$-fold repeated game $G^T$ has a unique subgame perfect equilibrium in which
$a^t = a^*$ for all $t \in \{1, \ldots, T\}$.
\end{theorem}

\begin{proof}
We proceed by backward induction on $t$.

\textbf{Base case ($t = T$):} In the final period, regardless of history $h^T$, the
subgame is identical to the one-shot stage game $G$. Since $a^*$ is the unique NE of $G$,
every SPE must play $a^*$ in period $T$.

\textbf{Inductive step:} Suppose that in every SPE, $a^{\tau} = a^*$ for all $\tau \in
\{t+1, \ldots, T\}$, regardless of the history at period $t+1$. Consider the subgame
beginning at an arbitrary history $h^t$. The continuation payoffs from period $t+1$ onward
are fixed at $\frac{1}{T}\sum_{\tau=t+1}^{T} g_i(a^*) = \frac{T-t}{T} g_i(a^*)$,
independent of the action chosen at period $t$ (since all future play is $a^*$ regardless
of history). Therefore, player $i$'s effective payoff at period $t$ is
\[
    \frac{1}{T} g_i(a^t) + \frac{T-t}{T} g_i(a^*),
\]
and maximizing this over $a_i^t$ is equivalent to maximizing $g_i(a_i^t, a_{-i}^t)$ in the
stage game. Since $a^*$ is the unique NE, $a^t = a^*$ in every SPE.

By induction, $a^t = a^*$ for all $t$, proving that the unique SPE is to play $a^*$ in
every period.
\end{proof}

\begin{warning}[The Unraveling Trap]
Theorem~\ref{thm:unraveling} shows that finite repetition of the Prisoner's Dilemma (which
has a unique NE at mutual defection) cannot sustain cooperation under subgame perfection.
This is counterintuitive and motivates the study of infinitely repeated games, where the
``last period'' argument cannot be applied.
\end{warning}

\subsection{Multiple Stage-Game Equilibria}

When the stage game has multiple Nash equilibria, repetition \emph{can} create new
equilibrium outcomes.

\begin{proposition}[New Equilibria Through Repetition]
\label{prop:multiple-ne-repetition}
If the stage game $G$ has at least two Nash equilibria yielding different payoffs for some
player, then the two-fold repeated game $G^2$ may have SPE outcomes in which first-period
play is not a stage-game NE.
\end{proposition}

\begin{proof}
We give a constructive proof. Let $G$ have two Nash equilibria $a^*$ and $\hat{a}$ with
$g_i(a^*) > g_i(\hat{a})$ for all $i$ (i.e., $a^*$ Pareto-dominates $\hat{a}$). Suppose
there exists an action profile $\tilde{a}$ (not necessarily a NE) with
$g_i(\tilde{a}) > g_i(a^*)$ for all $i$, and suppose the gain from deviating unilaterally
in the stage game satisfies
\[
    \max_{a_i' \in A_i} g_i(a_i', \tilde{a}_{-i}) - g_i(\tilde{a})
    < g_i(a^*) - g_i(\hat{a}).
\]
Consider the following strategy for each player $i$ in $G^2$:
\begin{itemize}
    \item Period 1: Play $\tilde{a}_i$.
    \item Period 2: If $a^1 = \tilde{a}$, play $a_i^*$. Otherwise, play $\hat{a}_i$.
\end{itemize}

In period 2, players play a stage-game NE ($a^*$ or $\hat{a}$) depending on period-1
history, so the strategy is sequentially rational in period 2.

In period 1, if player $i$ deviates to $a_i'$ while others play $\tilde{a}_{-i}$, the
payoff change is:
\begin{align*}
    &\frac{1}{2}\bigl[g_i(a_i', \tilde{a}_{-i}) + g_i(\hat{a})\bigr]
    - \frac{1}{2}\bigl[g_i(\tilde{a}) + g_i(a^*)\bigr] \\
    &= \frac{1}{2}\bigl[g_i(a_i', \tilde{a}_{-i}) - g_i(\tilde{a})\bigr]
    - \frac{1}{2}\bigl[g_i(a^*) - g_i(\hat{a})\bigr] < 0,
\end{align*}
by the assumed inequality. Thus no player wants to deviate, and the profile is an SPE
in which first-period play $\tilde{a}$ is \emph{not} a stage-game NE.
\end{proof}

\begin{example}[Coordination Game]
\label{ex:coordination-repeated}
Consider the two-player coordination game with payoffs:
\[
\begin{array}{c|cc}
  & L & R \\ \hline
T & 3,3 & 0,0 \\
B & 0,0 & 1,1
\end{array}
\]
This has two pure NE: $(T,L)$ with payoffs $(3,3)$ and $(B,R)$ with payoffs $(1,1)$.
In the two-fold repetition, the ``punishment'' equilibrium $(B,R)$ can incentivize
cooperation on outcomes that are not stage-game NE, provided the incentive condition is
satisfied.
\end{example}


%% ============================================================
%% SECTION 3: Infinitely Repeated Games
%% ============================================================
\section{Infinitely Repeated Games}
\label{sec:infinitely-repeated}

\subsection{Model}

To circumvent the unraveling argument, we consider infinite repetition with discounting.

\begin{definition}[Infinitely Repeated Game]
\label{def:infinitely-repeated}
Given a stage game $G$ and a discount factor $\delta \in (0,1)$, the \emph{infinitely
repeated game} $G^\infty(\delta)$ is defined by:
\begin{enumerate}[label=(\alph*)]
    \item In each period $t \in \{1, 2, 3, \ldots\}$, players simultaneously choose
    actions $a^t \in A$.
    \item Histories, strategies are defined as in the finite case, but with $t$ ranging
    over $\N$.
    \item Player $i$'s payoff under strategy profile $s$ is the \emph{normalized discounted
    sum}:
    \[
        u_i(s) = (1 - \delta) \sum_{t=1}^{\infty} \delta^{t-1} g_i(a^t).
    \]
\end{enumerate}
The normalization factor $(1-\delta)$ ensures that payoffs lie in the same range as
stage-game payoffs: if $a^t = a$ for all $t$, then $u_i = g_i(a)$.
\end{definition}

\subsection{Feasible and Individually Rational Payoffs}

\begin{definition}[Feasible Payoff Set]
\label{def:feasible-payoffs}
The \emph{feasible payoff set} of the stage game $G$ is
\[
    V = \conv\{g(a) : a \in A\} = \conv\{(g_1(a), \ldots, g_n(a)) : a \in A\} \subset \R^n,
\]
the convex hull of all achievable stage-game payoff vectors.
\end{definition}

\begin{remark}
In the infinitely repeated game, any payoff vector $v \in V$ is achievable by cycling
through action profiles with appropriate frequencies. For instance, if $v = \lambda g(a) +
(1-\lambda) g(a')$ for some $\lambda \in [0,1]$, then playing $a$ for a fraction $\lambda$
of periods and $a'$ for the remaining fraction achieves payoff $v$ in the limit as
$\delta \to 1$.
\end{remark}

\begin{definition}[Minimax Value and Individually Rational Payoffs]
\label{def:minimax-ir}
The \emph{minimax value} (or \emph{minmax value}) of player $i$ in pure strategies is
\[
    \underline{v}_i = \min_{a_{-i} \in A_{-i}} \max_{a_i \in A_i} g_i(a_i, a_{-i}).
\]
The \emph{mixed-strategy minimax value} is
\[
    \underline{v}_i^{\text{mix}} = \min_{\alpha_{-i} \in \Delta(A_{-i})} \max_{a_i \in A_i}
    \E_{\alpha_{-i}}[g_i(a_i, a_{-i})],
\]
where $\Delta(A_{-i}) = \prod_{j \neq i} \Delta(A_j)$ is the set of independent mixed
strategy profiles for $i$'s opponents. A payoff vector $v \in \R^n$ is
\emph{individually rational} if $v_i \geq \underline{v}_i^{\text{mix}}$ for all $i$.
\end{definition}

\begin{proposition}
\label{prop:ne-payoff-lower-bound}
In any Nash equilibrium of $G^\infty(\delta)$, each player $i$ receives payoff
$u_i \geq \underline{v}_i^{\text{mix}}$.
\end{proposition}

\begin{proof}
In any NE, player $i$ can deviate to the strategy of playing
$\argmax_{a_i} \E_{\alpha_{-i}^t}[g_i(a_i, a_{-i}^t)]$ in every period $t$ (where
$\alpha_{-i}^t$ is the marginal distribution of opponents' play at time $t$ induced by the
equilibrium). This guarantees a stage payoff of at least $\underline{v}_i^{\text{mix}}$ in
every period, hence a total payoff of at least $\underline{v}_i^{\text{mix}}$.
Since the NE strategy must yield payoff at least as high as any deviation, we get
$u_i \geq \underline{v}_i^{\text{mix}}$.
\end{proof}

\subsection{Trigger Strategies}

\begin{definition}[Grim Trigger Strategy]
\label{def:grim-trigger}
Given a target action profile $a^* \in A$, the \emph{grim trigger strategy} for player $i$
is:
\[
    s_i^t(h^t) = \begin{cases}
        a_i^* & \text{if } a^\tau = a^* \text{ for all } \tau < t, \\
        m_i   & \text{otherwise},
    \end{cases}
\]
where $m_i \in \argmin_{\alpha_i} \max_{a_j} g_j(a_j, \alpha_{-j})$ is player $i$'s
component of a minimax punishment profile against any deviator.
\end{definition}

\begin{example}[Cooperation in the Repeated Prisoner's Dilemma]
\label{ex:repeated-pd}
Consider the Prisoner's Dilemma with payoffs:
\[
\begin{array}{c|cc}
  & C & D \\ \hline
C & 3,3 & 0,4 \\
D & 4,0 & 1,1
\end{array}
\]
The unique stage-game NE is $(D,D)$ with payoffs $(1,1)$. The minimax value for each
player is $\underline{v}_i = 1$. Consider the grim trigger strategy profile that plays
$(C,C)$ and punishes with $(D,D)$ forever.

\textbf{Cooperation payoff:} If both cooperate forever, each gets $u_i = 3$.

\textbf{Deviation payoff:} If player $i$ defects in period $t$ (while the opponent
cooperates), she gets $4$ in that period and $1$ thereafter:
\[
    u_i^{\text{dev}} = (1-\delta)\bigl[\delta^{t-1} \cdot 4 + \sum_{\tau=t+1}^{\infty}
    \delta^{\tau-1} \cdot 1\bigr]
    + (1-\delta) \sum_{\tau=1}^{t-1} \delta^{\tau-1} \cdot 3.
\]
The best deviation is at $t=1$:
\[
    u_i^{\text{dev}} = (1-\delta) \cdot 4 + \delta \cdot 1 = 4 - 3\delta.
\]
Cooperation is sustained iff $3 \geq 4 - 3\delta$, i.e., $\delta \geq \frac{1}{3}$.
\end{example}

\begin{rlconnection}[Discount Factors and Episode Length in RL]
The discount factor $\delta$ in repeated games is analogous to the discount factor $\gamma$
in MDPs. In MARL, if $\gamma$ is too small (agents are too myopic), cooperative equilibria
may not be sustainable. This directly parallels the threshold $\delta \geq 1/3$ in the
repeated Prisoner's Dilemma: myopic agents (small $\gamma$) cannot learn to cooperate.
\end{rlconnection}

\begin{definition}[Tit-for-Tat]
\label{def:tit-for-tat}
In a two-player game, the \emph{tit-for-tat} strategy for player $i$ is:
\[
    s_i^1 = C, \qquad s_i^t(h^t) = a_j^{t-1} \quad \text{for } t \geq 2,
\]
where $j \neq i$. That is, start by cooperating, then mimic the opponent's previous action.
\end{definition}

\begin{remark}
Tit-for-tat is a ``forgiving'' strategy: unlike grim trigger, punishment is temporary.
Axelrod's famous tournaments showed that tit-for-tat performs well in evolutionary settings,
which connects to the multi-agent RL idea that simple, interpretable strategies can be
surprisingly effective.
\end{remark}


%% ============================================================
%% SECTION 4: Folk Theorems
%% ============================================================
\section{Folk Theorems}
\label{sec:folk-theorems}

The folk theorems characterize the equilibrium payoff set of infinitely repeated games.
Informally, ``almost anything'' can be sustained as an equilibrium if players are patient
enough.

\subsection{Nash Folk Theorem}

\begin{theorem}[Nash Folk Theorem]
\label{thm:nash-folk}
Let $G$ be a finite normal-form game and $V = \conv\{g(a) : a \in A\}$ the feasible payoff
set. For any payoff vector $v \in V$ satisfying $v_i > \underline{v}_i^{\text{mix}}$ for
all $i$ (strictly individually rational), there exists $\underline{\delta} \in (0,1)$ such
that for all $\delta \in (\underline{\delta}, 1)$, there is a Nash equilibrium of
$G^\infty(\delta)$ with payoff vector $v$.
\end{theorem}

\begin{proof}
We construct a Nash equilibrium yielding payoff $v$.

\textbf{Step 1: Achieving $v$ through action sequences.}
Since $v \in V = \conv\{g(a) : a \in A\}$, there exist action profiles
$a(1), a(2), \ldots, a(K) \in A$ and weights $\lambda_1, \ldots, \lambda_K \geq 0$ with
$\sum_k \lambda_k = 1$ such that $v = \sum_k \lambda_k g(a(k))$. We can approximate $v$
by cycling through these action profiles. Specifically, for any $\varepsilon > 0$ and
$\delta$ close enough to 1, we can construct a deterministic sequence of action profiles
$(a^t)_{t=1}^\infty$ such that the discounted average payoff
$(1-\delta)\sum_{t=1}^\infty \delta^{t-1} g(a^t)$ is within $\varepsilon$ of $v$.

More precisely, define a cycle of length $L$ where action profile $a(k)$ is played for
$\lceil \lambda_k L \rceil$ periods (adjusting the last to make the total exactly $L$).
As $\delta \to 1$, the discounted average of any periodic sequence converges to the
time-average, which converges to $v$ as $L \to \infty$. Fix $\varepsilon > 0$ small enough
that $v_i - \varepsilon > \underline{v}_i^{\text{mix}}$ for all $i$, and choose $L$ and
$\underline{\delta}_1$ accordingly.

\textbf{Step 2: Punishment strategies.}
For each player $i$, let $\alpha_{-i}^{m,i} \in \Delta(A_{-i})$ be a mixed minimax
punishment profile:
$\max_{a_i} \E_{\alpha_{-i}^{m,i}}[g_i(a_i, a_{-i})] = \underline{v}_i^{\text{mix}}$.
In the punishment phase, players $j \neq i$ play according to $\alpha_{-i}^{m,i}$
independently.

\textbf{Step 3: Strategy construction.}
Define the following strategy for each player $j$:
\begin{itemize}
    \item \emph{Cooperation phase:} Play the sequence $(a^t)_{t=1}^\infty$ from Step 1.
    \item \emph{Punishment phase:} If some player $i$ is the first to deviate from the
    prescribed sequence in some period $\hat{t}$, then for all $t > \hat{t}$, all players
    $j \neq i$ play according to $\alpha_{-i}^{m,i}$ (punishing $i$ forever).
\end{itemize}

\textbf{Step 4: Verifying the Nash equilibrium condition.}
Consider whether player $i$ wants to deviate. If player $i$ follows the prescribed strategy,
she receives payoff within $\varepsilon$ of $v_i$ (for $\delta > \underline{\delta}_1$).

If player $i$ deviates at any period $\hat{t}$, she receives at most
$\max_a g_i(a) \cdot (1-\delta)\delta^{\hat{t}-1}$ in that period, plus the punishment
continuation payoff. In the punishment phase, the opponents play the mixed minimax strategy,
so player $i$'s best response yields at most $\underline{v}_i^{\text{mix}}$ per period.
Thus the deviation payoff is at most:
\[
    u_i^{\text{dev}} \leq (1-\delta) \sum_{\tau=1}^{\hat{t}} \delta^{\tau-1} M
    + \delta^{\hat{t}} \underline{v}_i^{\text{mix}}
\]
where $M = \max_a |g_i(a)|$.

As $\delta \to 1$, the first term (a finite number of periods) vanishes and
$u_i^{\text{dev}} \to \underline{v}_i^{\text{mix}} < v_i - \varepsilon$. Therefore,
for $\delta$ large enough (say $\delta > \underline{\delta}_2$), deviation is unprofitable.

Taking $\underline{\delta} = \max(\underline{\delta}_1, \underline{\delta}_2)$ completes
the proof.

\textbf{Note on Nash vs.\ SPE:} This equilibrium is Nash but not necessarily subgame
perfect: punishing players have no incentive to carry out the punishment (they are
indifferent or may prefer not to punish). The subgame perfect folk theorem below addresses
this.
\end{proof}

\begin{keyresult}[Nash Folk Theorem -- Interpretation]
Any feasible, strictly individually rational payoff vector is achievable as a Nash
equilibrium payoff for sufficiently patient players. This means the set of NE payoffs
of the repeated game ``fills up'' the feasible individually rational region as
$\delta \to 1$.
\end{keyresult}

\subsection{Subgame Perfect Folk Theorem}

The Nash folk theorem has a weakness: punishment threats may not be credible. The punishers
might prefer to deviate from punishment. The subgame perfect folk theorem (Fudenberg and
Maskin, 1986) resolves this by constructing self-enforcing punishment schemes.

\begin{definition}[Full Dimensionality Condition]
\label{def:full-dim}
The stage game $G$ satisfies the \emph{full dimensionality condition} if the feasible payoff
set $V = \conv\{g(a) : a \in A\}$ has dimension $n$ (the number of players), i.e.,
$V$ has nonempty interior in $\R^n$.
\end{definition}

\begin{theorem}[Subgame Perfect Folk Theorem (Fudenberg--Maskin)]
\label{thm:spe-folk}
Let $G$ be a finite $n$-player game satisfying the full dimensionality condition. For any
feasible payoff vector $v \in V$ with $v_i > \underline{v}_i^{\text{mix}}$ for all $i$,
there exists $\underline{\delta} \in (0,1)$ such that for all
$\delta \in (\underline{\delta}, 1)$, there is a subgame perfect equilibrium of
$G^\infty(\delta)$ with payoff vector $v$.
\end{theorem}

\begin{proof}
The proof is constructive. We build a strategy profile with multiple phases: a cooperation
phase and player-specific punishment phases, arranged so that each phase is self-enforcing.

\textbf{Step 1: Setup.}
Fix $v \in V$ with $v_i > \underline{v}_i^{\text{mix}}$ for all $i$. Since $V$ has full
dimension, for each player $i$ there exists a feasible payoff vector $w^i \in V$ with:
\begin{enumerate}[label=(\roman*)]
    \item $w^i_j > \underline{v}_j^{\text{mix}}$ for all $j \neq i$ (the punishers receive
    strictly above their minimax values during the punishment of $i$),
    \item $w^i_i < v_i$ (player $i$'s payoff during her own punishment phase is below her
    cooperative payoff; we can choose $w^i_i$ close to $\underline{v}_i^{\text{mix}}$).
\end{enumerate}
The full dimensionality condition ensures such $w^i$ exist: since $V$ has nonempty interior,
we can find vectors that independently adjust each player's payoff.

Let $m^i_{-i} \in \Delta(A_{-i})$ be a mixed minimax strategy profile against player $i$.

\textbf{Step 2: Strategy construction (three-phase strategy).}

\emph{Phase 0 (Cooperation):} Play the action sequence that yields payoff $v$ (as in the
Nash folk theorem proof).

\emph{Phase I$_i$ (Punishment of player $i$):} Lasts for $L$ periods. During this phase,
the opponents of $i$ play the minimax strategy $m^i_{-i}$ in each period, and player $i$
plays any best response (which yields at most $\underline{v}_i^{\text{mix}}$ per period).

\emph{Phase II$_i$ (Reward after punishment):} After the $L$-period punishment phase, play
switches to the action sequence yielding payoff $w^i$ forever. This rewards the punishers
(who received above their minimax during punishment) and keeps $i$'s continuation payoff
below $v_i$ if $L$ is large enough.

\emph{Transitions:}
\begin{itemize}
    \item If player $j$ is the first to deviate in Phase 0, switch to Phase I$_j$.
    \item If player $j$ deviates during Phase I$_i$ (and $j \neq i$), restart Phase I$_j$
    (punish the deviating punisher).
    \item If player $i$ deviates during her own punishment Phase I$_i$, restart Phase I$_i$.
\end{itemize}

\textbf{Step 3: Choosing $L$.}
We need $L$ large enough that deviation is unprofitable. Let $M = \max_{a,i} |g_i(a)|$.
During Phase I$_i$, player $i$ gets at most $\underline{v}_i^{\text{mix}}$ per period.
After Phase I$_i$, player $i$ gets $w^i_i$ per period. Thus player $i$'s continuation payoff
from Phase I$_i$ is:
\[
    (1-\delta) \sum_{t=1}^{L} \delta^{t-1} \underline{v}_i^{\text{mix}}
    + \delta^L w^i_i
    = (1 - \delta^L) \underline{v}_i^{\text{mix}} + \delta^L w^i_i.
\]
We need this to be less than $v_i$ (so punishment deters deviation from Phase 0) and for
the punishers' payoffs to make Phase I self-enforcing.

\textbf{Step 4: Verifying sequential rationality.}

\emph{Phase 0 -- Player $i$ considers deviating:}
Deviation yields at most $M$ in one period, then Phase I$_i$ begins. Compliance yields
$v_i$ forever. The deviation is unprofitable when:
\[
    v_i \geq (1-\delta) M + \delta \bigl[(1 - \delta^L)
    \underline{v}_i^{\text{mix}} + \delta^L w^i_i\bigr].
\]
As $\delta \to 1$, the right side approaches $(1-\delta)M + \delta[(1-\delta^L)
\underline{v}_i^{\text{mix}} + \delta^L w^i_i]$. For $\delta$ close to 1 and $L$ large,
$(1-\delta)M \to 0$ and the continuation approaches a weighted average of
$\underline{v}_i^{\text{mix}}$ and $w^i_i$, both strictly below $v_i$ (since $w^i_i < v_i$
and $\underline{v}_i^{\text{mix}} < v_i$). So the inequality holds for $\delta$ large
enough.

\emph{Phase I$_i$ -- Punisher $j \neq i$ considers deviating:}
If $j$ deviates during Phase I$_i$, Phase I$_j$ begins (punishing $j$). Compliance with
Phase I$_i$ yields: the minimax strategy payoff for $L-t$ remaining periods (at worst
$-M$ per period but bounded), followed by $w^i_j > \underline{v}_j^{\text{mix}}$ forever.
Deviation yields at most $M$ for one period, then Phase I$_j$ punishment.

For $\delta$ close to 1, the continuation payoff from compliance converges to $w^i_j >
\underline{v}_j^{\text{mix}}$, while deviation leads eventually to $w^j_j$ which can be
chosen close to $\underline{v}_j^{\text{mix}}$. So compliance dominates for large $\delta$.

\emph{Phase I$_i$ -- Player $i$ considers deviating:}
Deviating restarts Phase I$_i$, gaining at most $M$ in one period but adding $L$ more
punishment periods. Compliance ends the punishment sooner and leads to $w^i_i$ forever.
For large $\delta$, the cost of extending punishment exceeds any one-period gain.

Since all deviations are unprofitable for $\delta > \underline{\delta}$ (where
$\underline{\delta}$ depends on the game parameters $M, v, w^i,
\underline{v}_i^{\text{mix}}, L$), the constructed strategy profile is a subgame perfect
equilibrium.
\end{proof}

\begin{remark}
The full dimensionality condition is needed only for $n \geq 3$ players: it allows
constructing punishment payoffs $w^i$ that reward the punishers while punishing the
deviator. For two-player games, Fudenberg and Maskin (1986) showed the folk theorem holds
without this condition, since minmaxing the opponent is the only thing the punisher can do.
\end{remark}

\begin{keyresult}[Subgame Perfect Folk Theorem -- Interpretation]
For patient enough players (large $\delta$), the set of SPE payoffs of the infinitely
repeated game is exactly the set of feasible, strictly individually rational payoff
vectors (under the full dimensionality condition). This is a much stronger result than the
Nash folk theorem because SPE requires credible (self-enforcing) punishments.
\end{keyresult}

\begin{rlconnection}[Folk Theorems and Multi-Agent RL]
The folk theorems have profound implications for MARL:
\begin{enumerate}[label=(\roman*)]
    \item \textbf{Multiplicity of equilibria:} In repeated general-sum games, there are
    typically \emph{many} equilibria. This means that independent learners may converge to
    vastly different outcomes depending on initial conditions, learning rates, and
    exploration schedules.
    \item \textbf{Emergent cooperation:} The folk theorems show that cooperation can be
    sustained by punishment threats, explaining why self-play agents sometimes learn
    cooperative strategies in iterated social dilemmas.
    \item \textbf{Reward shaping:} The discount factor $\delta$ (analogous to $\gamma$ in
    RL) determines the set of achievable equilibria. Curriculum designers can influence
    emergent behavior by adjusting $\gamma$.
\end{enumerate}
\end{rlconnection}


%% ============================================================
%% SECTION 5: Zero-Sum Games -- Deeper Treatment
%% ============================================================
\section{Zero-Sum Games: Fictitious Play and No-Regret Learning}
\label{sec:zero-sum}

We return to two-player zero-sum games $G = (\{1,2\}, A_1, A_2, g)$ where
$g_1(a) = -g_2(a) = g(a)$. Recall from Lesson 11 the minimax theorem:
$\max_{\alpha_1} \min_{\alpha_2} \E[g(\alpha_1, \alpha_2)]
= \min_{\alpha_2} \max_{\alpha_1} \E[g(\alpha_1, \alpha_2)] = v^*$ (the value of the game).

\subsection{Fictitious Play}

\begin{definition}[Fictitious Play]
\label{def:fictitious-play}
\emph{Fictitious play} is a learning dynamic where each player best-responds to the
empirical frequency of the opponent's past actions. Formally, let $a_j^1, a_j^2, \ldots$
be the actions played by player $j$. The empirical frequency after $T$ rounds is
\[
    \hat{\alpha}_j^T(a_j) = \frac{1}{T} \sum_{t=1}^T \indicator\{a_j^t = a_j\}
    \qquad \text{for each } a_j \in A_j.
\]
At round $T+1$, player $i$ plays $a_i^{T+1} \in \BR_i(\hat{\alpha}_j^T)$, a best response
to the empirical distribution of the opponent.
\end{definition}

\begin{theorem}[Convergence of Fictitious Play in $2 \times 2$ Zero-Sum Games]
\label{thm:fp-2x2}
In any $2 \times 2$ zero-sum game, the empirical frequencies
$(\hat{\alpha}_1^T, \hat{\alpha}_2^T)$ under fictitious play converge to a Nash equilibrium
as $T \to \infty$.
\end{theorem}

\begin{proof}
Let the payoff matrix be
\[
    G = \begin{pmatrix} a & b \\ c & d \end{pmatrix}
\]
with player 1 choosing rows and player 2 choosing columns, so $g_1 = G$, $g_2 = -G$.

Consider the function $\Phi^T = \max_{a_1} \sum_{t=1}^T g(a_1, a_2^t)$. This is player 1's
``fictitious play payoff''---the best she could have done against the sequence of player 2's
actions if constrained to a fixed action.

We analyze how $\Phi^T$ evolves. Note that $\Phi^T = \max\{R_1^T, R_2^T\}$ where
$R_k^T = \sum_{t=1}^T g(k, a_2^t)$ is the cumulative payoff of row $k$.

In a $2 \times 2$ game, the best response of player 1 to $\hat{\alpha}_2^T$ is row 1 if
$R_1^T > R_2^T$ and row 2 if $R_1^T < R_2^T$ (with either at ties). Define
$D^T = R_1^T - R_2^T$.

\emph{Claim:} $|D^T|$ is bounded. To see this, note that when $D^T > 0$, player 1 plays
row 1, and player 2 best-responds to the empirical frequency $\hat{\alpha}_1^T$ which puts
more weight on row 1. Player 2's best response tends to be the column that makes row 1
less attractive, which decreases $D^{T+1}$. Similarly when $D^T < 0$. The key observation
in the $2 \times 2$ case is that $D^T$ oscillates around 0 with bounded amplitude
(bounded by $\max |g(a)|$), because each step changes $D^T$ by at most
$\max_{a_2} |g(1,a_2) - g(2,a_2)|$.

More precisely, define $V_1^T = \frac{1}{T}\Phi^T$ (player 1's ``best fixed action''
average payoff). We have
\[
    V_1^T = \frac{1}{T} \max_{a_1} \sum_{t=1}^T g(a_1, a_2^t)
    \geq \frac{1}{T} \sum_{t=1}^T g(a_1^t, a_2^t),
\]
since the best fixed action does at least as well as the actual sequence.

Similarly, define $V_2^T = \frac{1}{T} \max_{a_2} \sum_{t=1}^T (-g(a_1^t, a_2))$ for
player 2. By the minimax theorem for the stage game:
\[
    V_1^T \geq v^* \geq -V_2^T.
\]
Furthermore, $V_1^T - v^* \to 0$ and $V_2^T + v^* \to 0$ because the ``regret''
$V_1^T - \frac{1}{T}\sum_t g(a_1^t, a_2^t)$ can be bounded by the bounded oscillation of
$D^T/T \to 0$.

Since $|D^T|$ is bounded, $\frac{D^T}{T} \to 0$, which means $\frac{R_1^T}{T}$ and
$\frac{R_2^T}{T}$ converge to the same limit. This implies that $\hat{\alpha}_2^T$
converges to making player 1 indifferent between rows, which is exactly the NE condition
in a $2 \times 2$ game with a completely mixed equilibrium. If the game has a pure
equilibrium, fictitious play converges to it in finitely many steps.

Thus $(\hat{\alpha}_1^T, \hat{\alpha}_2^T) \to (\alpha_1^*, \alpha_2^*)$, a Nash
equilibrium.
\end{proof}

\begin{theorem}[Robinson's Theorem: Fictitious Play in General Zero-Sum Games]
\label{thm:robinson}
In any finite two-player zero-sum game, the empirical frequencies under fictitious play
converge to the set of Nash equilibria.
\end{theorem}

\begin{proof}[Proof sketch]
The full proof (Robinson, 1951) uses a potential function argument. Define
\[
    \Phi_1^T = \max_{a_1 \in A_1} \sum_{t=1}^T g(a_1, a_2^t), \qquad
    \Phi_2^T = \max_{a_2 \in A_2} \sum_{t=1}^T (-g(a_1^t, a_2)).
\]
Then $\frac{1}{T}\Phi_1^T \geq v^*$ and $\frac{1}{T}\Phi_2^T \geq -v^*$ (each player's
best fixed action does at least as well as the game value against the empirical opponent
distribution).

The key step is showing that $\frac{1}{T}(\Phi_1^T + \Phi_2^T) \to 0$. This follows
because $\Phi_1^T + \Phi_2^T$ grows sublinearly: when player 1 plays a best response to
$\hat{\alpha}_2^T$, the increase in $\Phi_1^{T+1} - \Phi_1^T$ is exactly
$g(a_1^{T+1}, a_2^{T+1})$, while the increase $\Phi_2^{T+1} - \Phi_2^T \leq
-g(a_1^{T+1}, a_2^{T+1}) + C$ for a constant $C$ related to the number of actions.
Carefully tracking these increments shows $\Phi_1^T + \Phi_2^T = O(\sqrt{T})$.

Since $\frac{1}{T}\Phi_1^T \to v^*$ and $\frac{1}{T}\Phi_2^T \to -v^*$, any limit point
$(\alpha_1^*, \alpha_2^*)$ of $(\hat{\alpha}_1^T, \hat{\alpha}_2^T)$ satisfies the NE
conditions: $\alpha_1^*$ is a best response to $\alpha_2^*$ and vice versa.
\end{proof}

\begin{rlconnection}[Fictitious Play and Policy-Space Response Oracles (PSRO)]
Fictitious play is a precursor to the \emph{PSRO} framework in modern MARL. In PSRO,
each player maintains a population of policies (analogous to the empirical frequency in FP),
and computes a best response to the current mixture---exactly as in fictitious play, but
with neural-network policies instead of tabular strategies. The convergence guarantees of
FP in zero-sum games provide theoretical backing for PSRO's effectiveness in competitive
MARL.
\end{rlconnection}

\subsection{No-Regret Learning}

\begin{definition}[External Regret]
\label{def:external-regret}
Player $i$'s \emph{external regret} after $T$ rounds is
\[
    R_i^T = \max_{a_i \in A_i} \sum_{t=1}^T g_i(a_i, a_{-i}^t)
    - \sum_{t=1}^T g_i(a_i^t, a_{-i}^t).
\]
A learning algorithm has \emph{no external regret} (or is \emph{Hannan consistent}) if
$R_i^T / T \to 0$ as $T \to \infty$.
\end{definition}

\begin{definition}[Multiplicative Weights Update (MWU)]
\label{def:mwu}
Let $\eta > 0$ be a learning rate. Player $i$ maintains weights $w_i^t(a_i)$ for each
action $a_i \in A_i$:
\begin{align*}
    w_i^1(a_i) &= 1 \quad \text{for all } a_i, \\
    w_i^{t+1}(a_i) &= w_i^t(a_i) \cdot \exp\bigl(\eta \cdot g_i(a_i, a_{-i}^t)\bigr).
\end{align*}
The mixed strategy at time $t$ is $\alpha_i^t(a_i) = w_i^t(a_i) / \sum_{a_i'} w_i^t(a_i')$.
\end{definition}

\begin{theorem}[MWU Regret Bound]
\label{thm:mwu-regret}
If $|g_i(a)| \leq 1$ for all $a$, then with $\eta = \sqrt{\frac{\ln |A_i|}{T}}$, the
multiplicative weights update achieves
\[
    R_i^T \leq 2\sqrt{T \ln |A_i|}.
\]
In particular, $R_i^T / T \to 0$, so MWU is a no-regret algorithm.
\end{theorem}

\begin{proof}
Define the potential $W^t = \sum_{a_i \in A_i} w_i^t(a_i)$, so $W^1 = |A_i|$.

\textbf{Upper bound on $\ln(W^{T+1}/W^1)$:}
\begin{align*}
    \frac{W^{t+1}}{W^t}
    &= \frac{\sum_{a_i} w_i^t(a_i) \exp(\eta \cdot g_i(a_i, a_{-i}^t))}{W^t} \\
    &= \sum_{a_i} \alpha_i^t(a_i) \exp(\eta \cdot g_i(a_i, a_{-i}^t)) \\
    &\leq \sum_{a_i} \alpha_i^t(a_i) \bigl(1 + \eta g_i(a_i, a_{-i}^t)
    + \eta^2 g_i(a_i, a_{-i}^t)^2\bigr),
\end{align*}
using $e^x \leq 1 + x + x^2$ for $|x| \leq 1$ (valid when $\eta \leq 1$ and
$|g_i| \leq 1$). Therefore:
\begin{align*}
    \frac{W^{t+1}}{W^t}
    &\leq 1 + \eta \sum_{a_i} \alpha_i^t(a_i) g_i(a_i, a_{-i}^t)
    + \eta^2 \\
    &\leq \exp\bigl(\eta \cdot g_i(\alpha_i^t, a_{-i}^t) + \eta^2\bigr),
\end{align*}
where $g_i(\alpha_i^t, a_{-i}^t) = \E_{a_i \sim \alpha_i^t}[g_i(a_i, a_{-i}^t)]$ and
we used $1 + x \leq e^x$. Telescoping:
\[
    \ln \frac{W^{T+1}}{W^1} \leq \eta \sum_{t=1}^T g_i(\alpha_i^t, a_{-i}^t) + \eta^2 T.
\]

\textbf{Lower bound on $\ln(W^{T+1}/W^1)$:}
For any fixed action $a_i^*$:
\[
    W^{T+1} \geq w_i^{T+1}(a_i^*) = \exp\Bigl(\eta \sum_{t=1}^T g_i(a_i^*, a_{-i}^t)\Bigr).
\]
Therefore:
\[
    \ln \frac{W^{T+1}}{W^1} \geq \eta \sum_{t=1}^T g_i(a_i^*, a_{-i}^t) - \ln |A_i|.
\]

\textbf{Combining:}
\[
    \eta \sum_{t=1}^T g_i(a_i^*, a_{-i}^t) - \ln |A_i|
    \leq \eta \sum_{t=1}^T g_i(\alpha_i^t, a_{-i}^t) + \eta^2 T.
\]
Rearranging:
\[
    \sum_{t=1}^T g_i(a_i^*, a_{-i}^t) - \sum_{t=1}^T g_i(\alpha_i^t, a_{-i}^t)
    \leq \frac{\ln |A_i|}{\eta} + \eta T.
\]
Since this holds for all $a_i^*$, taking the max and substituting
$\eta = \sqrt{\frac{\ln |A_i|}{T}}$:
\[
    R_i^T \leq \sqrt{T \ln |A_i|} + \sqrt{T \ln |A_i|} = 2\sqrt{T \ln |A_i|}.
    \qedhere
\]
\end{proof}

\begin{theorem}[No-Regret Implies NE in Zero-Sum Games]
\label{thm:no-regret-ne}
In a two-player zero-sum game, if both players use no-regret learning algorithms, then
the time-averaged strategy profile $(\bar{\alpha}_1^T, \bar{\alpha}_2^T)$ converges to the
set of Nash equilibria as $T \to \infty$, where
$\bar{\alpha}_i^T = \frac{1}{T}\sum_{t=1}^T \alpha_i^t$.
\end{theorem}

\begin{proof}
Let $g = g_1 = -g_2$. The no-regret condition for player 1 gives:
\[
    \max_{a_1} \frac{1}{T} \sum_{t=1}^T g(a_1, \alpha_2^t)
    \leq \frac{1}{T} \sum_{t=1}^T g(\alpha_1^t, \alpha_2^t) + \frac{R_1^T}{T}.
\]
By linearity of $g$ in the first argument and Jensen's inequality (or directly by linearity):
\[
    \max_{a_1} g(a_1, \bar{\alpha}_2^T)
    = \max_{a_1} \frac{1}{T} \sum_{t=1}^T g(a_1, \alpha_2^t)
    \leq \frac{1}{T} \sum_{t=1}^T g(\alpha_1^t, \alpha_2^t) + \varepsilon_1^T,
\]
where $\varepsilon_1^T = R_1^T / T \to 0$.

Similarly, the no-regret condition for player 2 (who minimizes $g$) gives:
\[
    \min_{a_2} g(\bar{\alpha}_1^T, a_2)
    = \min_{a_2} \frac{1}{T} \sum_{t=1}^T g(\alpha_1^t, a_2)
    \geq \frac{1}{T} \sum_{t=1}^T g(\alpha_1^t, \alpha_2^t) - \varepsilon_2^T,
\]
where $\varepsilon_2^T \to 0$.

Combining:
\[
    \min_{a_2} g(\bar{\alpha}_1^T, a_2)
    \geq \frac{1}{T}\sum_t g(\alpha_1^t, \alpha_2^t) - \varepsilon_2^T
    \geq \max_{a_1} g(a_1, \bar{\alpha}_2^T) - \varepsilon_1^T - \varepsilon_2^T.
\]

But by the minimax inequality (which always holds):
\[
    \max_{a_1} g(a_1, \bar{\alpha}_2^T)
    \geq \min_{a_2} g(\bar{\alpha}_1^T, a_2).
\]

Therefore:
\[
    0 \leq \max_{a_1} g(a_1, \bar{\alpha}_2^T) - \min_{a_2} g(\bar{\alpha}_1^T, a_2)
    \leq \varepsilon_1^T + \varepsilon_2^T \to 0.
\]

This means the ``duality gap'' converges to zero. Since
$\max_{a_1} g(a_1, \bar{\alpha}_2^T) \to v^*$ and
$\min_{a_2} g(\bar{\alpha}_1^T, a_2) \to v^*$, any limit point
$(\alpha_1^*, \alpha_2^*)$ of $(\bar{\alpha}_1^T, \bar{\alpha}_2^T)$ satisfies
$\max_{a_1} g(a_1, \alpha_2^*) = \min_{a_2} g(\alpha_1^*, a_2) = v^*$, which is the
Nash equilibrium condition.
\end{proof}

\begin{rlconnection}[No-Regret Learning in Multi-Agent RL]
No-regret algorithms are fundamental to competitive MARL. The convergence guarantee of
Theorem~\ref{thm:no-regret-ne} underpins algorithms like:
\begin{itemize}
    \item \textbf{Counterfactual Regret Minimization (CFR)} for extensive-form games
    (poker AI).
    \item \textbf{Online mirror descent} in policy space for continuous games.
    \item \textbf{V-learning} and other model-free algorithms that achieve sublinear regret
    against adversarial opponents.
\end{itemize}
The $O(\sqrt{T})$ regret rate of MWU translates to $O(1/\sqrt{T})$ convergence to NE in
zero-sum games, which can be accelerated by optimistic variants.
\end{rlconnection}


%% ============================================================
%% SECTION 6: Bayesian Games
%% ============================================================
\section{Bayesian Games (Incomplete Information)}
\label{sec:bayesian-games}

\subsection{Motivation and Model}

In many strategic settings, players are uncertain about others' payoff functions. A firm
does not know a competitor's cost structure; an agent does not know an opponent's reward
function. Harsanyi's (1967--68) framework models this via \emph{types}.

\begin{definition}[Bayesian Game]
\label{def:bayesian-game}
A \emph{Bayesian game} (or game of incomplete information) is a tuple
$\Gamma = (N, (A_i)_{i \in N}, (\Theta_i)_{i \in N}, p, (u_i)_{i \in N})$ where:
\begin{enumerate}[label=(\alph*)]
    \item $N = [n]$ is the player set.
    \item $A_i$ is the finite action set of player $i$.
    \item $\Theta_i$ is the finite \emph{type space} of player $i$. A type
    $\theta_i \in \Theta_i$ encodes player $i$'s private information.
    \item $p \in \Delta(\Theta)$ is the \emph{common prior} distribution over the type
    profile $\theta = (\theta_1, \ldots, \theta_n) \in \Theta = \prod_i \Theta_i$.
    \item $u_i : A \times \Theta \to \R$ is player $i$'s payoff function, depending on both
    the action profile and the type profile.
\end{enumerate}
The timing is: Nature draws $\theta \sim p$; each player $i$ observes only her own type
$\theta_i$; players simultaneously choose actions.
\end{definition}

\begin{definition}[Bayesian Nash Equilibrium]
\label{def:bne}
A \emph{(pure) strategy} in a Bayesian game is a function $s_i : \Theta_i \to A_i$ mapping
types to actions. A strategy profile $s^* = (s_1^*, \ldots, s_n^*)$ is a \emph{Bayesian
Nash equilibrium} (BNE) if, for every player $i$ and every type $\theta_i \in \Theta_i$:
\[
    s_i^*(\theta_i) \in \argmax_{a_i \in A_i}
    \E_{\theta_{-i} \sim p(\cdot | \theta_i)}\bigl[
    u_i\bigl(a_i, s_{-i}^*(\theta_{-i}), \theta_i, \theta_{-i}\bigr)\bigr],
\]
where $p(\theta_{-i} | \theta_i) = p(\theta_i, \theta_{-i}) / p(\theta_i)$ is the
conditional distribution of others' types given $\theta_i$.
\end{definition}

\begin{remark}
The BNE condition says: each type of each player best-responds to the expected play of
opponents, where the expectation is over opponents' types (which determine their actions
through the strategy profile).
\end{remark}

\subsection{Existence of BNE}

\begin{theorem}[Existence of BNE]
\label{thm:bne-existence}
Every finite Bayesian game (finite player set, finite action sets, finite type spaces) has
a Bayesian Nash equilibrium in mixed strategies.
\end{theorem}

\begin{proof}
Transform the Bayesian game into the ``agent normal form.'' For each player-type pair
$(i, \theta_i)$, create an ``agent'' who chooses action $a_i \in A_i$. This yields a
finite normal-form game with $\sum_i |\Theta_i|$ agents, each with action set $A_i$.
The payoff to agent $(i, \theta_i)$ when the action profile is
$(a_{j,\theta_j})_{j,\theta_j}$ is:
\[
    \sum_{\theta_{-i}} p(\theta_{-i} | \theta_i) \cdot
    u_i(a_{i,\theta_i}, (a_{j,\theta_j})_{j \neq i}, \theta_i, \theta_{-i}).
\]

This is a finite normal-form game, so by Nash's theorem (from Lesson 11), it has a Nash
equilibrium in mixed strategies. Any NE of the agent normal form corresponds to a BNE of
the original Bayesian game: the mixed strategy of agent $(i, \theta_i)$ specifies
$s_i(\theta_i) \in \Delta(A_i)$, and the NE condition for each agent is exactly the BNE
condition for type $\theta_i$ of player $i$.
\end{proof}

\subsection{Examples}

\begin{example}[Cournot Duopoly with Unknown Costs]
\label{ex:cournot-bayesian}
Two firms produce a homogeneous good with inverse demand $P(Q) = a - Q$ where
$Q = q_1 + q_2$. Firm $i$'s cost is $c_i q_i$. Suppose $c_1$ is known to both firms, but
$c_2$ is private information of firm 2: $c_2 \in \{c_L, c_H\}$ with
$\Prob(c_2 = c_L) = \pi$.

This is a Bayesian game with $\Theta_2 = \{c_L, c_H\}$, $\Theta_1 = \{c_1\}$ (trivial
type space for firm 1). Actions are quantities $q_i \geq 0$.

Firm 2's profit when type is $\theta_2 = c$: $u_2 = (a - q_1 - q_2 - c)q_2$.

In a BNE, firm 2 uses strategy $s_2 : \{c_L, c_H\} \to \R_+$ and firm 1 chooses $q_1$.
The first-order conditions give:
\begin{align*}
    q_1 &= \frac{a - c_1 - \pi s_2(c_L) - (1-\pi)s_2(c_H)}{2}, \\
    s_2(c_L) &= \frac{a - c_L - q_1}{2}, \\
    s_2(c_H) &= \frac{a - c_H - q_1}{2}.
\end{align*}
Solving this system of three equations yields the unique BNE. The key insight is that
firm 1's optimal quantity depends on the \emph{expected} cost of firm 2, and firm 2's
strategy conditions on its private type.
\end{example}

\begin{example}[First-Price Sealed-Bid Auction]
\label{ex:first-price-auction}
Two bidders have private values $v_i \sim \text{Uniform}[0,1]$ independently. Each bidder
submits a sealed bid $b_i$; the highest bidder wins and pays her bid. Payoff:
$u_i = (v_i - b_i)\indicator\{b_i > b_j\}$.

This is a Bayesian game with $\Theta_i = [0,1]$ (continuous types). We seek a symmetric
BNE where $b_i = \beta(v_i)$ for some increasing function $\beta$.

Given opponent's strategy $\beta$, bidder $i$ with value $v_i$ solves:
\[
    \max_{b} (v_i - b) \cdot \Prob(\beta(v_j) < b)
    = (v_i - b) \cdot \beta^{-1}(b).
\]
The first-order condition (substituting $b = \beta(v_i)$) and the boundary condition
$\beta(0) = 0$ yield $\beta(v_i) = v_i / 2$: each bidder bids half her value.

\textbf{Verification:} If bidder $j$ bids $v_j/2$, then bidder $i$'s problem is
$\max_b (v_i - b) \cdot 2b = \max_b 2bv_i - 2b^2$. The FOC gives $2v_i - 4b = 0$, so
$b = v_i/2$, confirming the BNE.
\end{example}

\begin{rlconnection}[Bayesian Games and Opponent Modeling in MARL]
In MARL, agents typically do not know each other's reward functions---precisely the
setting of Bayesian games. The Bayesian framework formalizes \emph{opponent modeling}:
\begin{itemize}
    \item The type $\theta_j$ can encode opponent $j$'s policy parameters or reward weights.
    \item The prior $p(\theta_{-i})$ represents agent $i$'s beliefs about opponents.
    \item Bayesian updating (observing opponents' actions and updating beliefs via Bayes'
    rule) is the principled approach to learning about opponents.
\end{itemize}
Algorithms like LOLA (Learning with Opponent-Learning Awareness) and Bayesian policy
reuse implement variants of this Bayesian framework.
\end{rlconnection}


%% ============================================================
%% SECTION 7: Perfect Bayesian Equilibrium
%% ============================================================
\section{Perfect Bayesian Equilibrium}
\label{sec:pbe}

\subsection{Motivation}

In dynamic games of incomplete information, the BNE concept is too weak: it does not
constrain behavior at off-equilibrium information sets. We need a refinement that combines
sequential rationality with belief consistency.

\begin{definition}[System of Beliefs]
\label{def:beliefs}
In an extensive-form game with incomplete information, a \emph{system of beliefs}
$\mu$ assigns to each information set $h$ a probability distribution $\mu(\cdot | h)$
over the nodes in that information set.
\end{definition}

\begin{definition}[Perfect Bayesian Equilibrium]
\label{def:pbe}
A \emph{perfect Bayesian equilibrium} (PBE) is a pair $(\sigma, \mu)$ of a behavioral
strategy profile $\sigma$ and a belief system $\mu$ such that:
\begin{enumerate}[label=(\alph*)]
    \item \textbf{Sequential rationality:} At every information set $h$ of every player $i$,
    the strategy $\sigma_i$ is optimal given beliefs $\mu(\cdot | h)$ and the continuation
    strategies $\sigma_{-i}$.
    \item \textbf{Belief consistency:} At every information set $h$ that is reached with
    positive probability under $\sigma$, the beliefs $\mu(\cdot | h)$ are derived from
    $\sigma$ using Bayes' rule:
    \[
        \mu(x | h) = \frac{\Prob(x \text{ is reached} | \sigma)}
        {\sum_{x' \in h} \Prob(x' \text{ is reached} | \sigma)}
        \quad \text{for all } x \in h.
    \]
    \item At information sets reached with zero probability (off the equilibrium path),
    beliefs are unrestricted (any distribution over nodes in the information set is allowed).
\end{enumerate}
\end{definition}

\subsection{Signaling Games}

A particularly important class of dynamic Bayesian games is \emph{signaling games}.

\begin{definition}[Signaling Game]
\label{def:signaling-game}
A \emph{signaling game} has two players:
\begin{itemize}
    \item \textbf{Sender (S):} Has private type $\theta \in \Theta$, drawn from prior
    $p(\theta)$. Observes $\theta$, then sends a message $m \in M$.
    \item \textbf{Receiver (R):} Observes $m$ (but not $\theta$), then takes action
    $a \in A$.
\end{itemize}
Payoffs $u_S(\theta, m, a)$ and $u_R(\theta, m, a)$ depend on all three: type, message,
and action.
\end{definition}

\begin{definition}[Equilibrium Types in Signaling Games]
\label{def:equilibrium-types}
\begin{enumerate}[label=(\alph*)]
    \item A PBE is \emph{separating} if different types send different messages:
    $\sigma_S(\theta) \neq \sigma_S(\theta')$ for $\theta \neq \theta'$.
    \item A PBE is \emph{pooling} if all types send the same message:
    $\sigma_S(\theta) = m^*$ for all $\theta$.
    \item A PBE is \emph{semi-separating} (or partially pooling) if some types pool and
    others separate.
\end{enumerate}
\end{definition}

\begin{example}[Spence's Job Market Signaling]
\label{ex:spence}
A worker has ability $\theta \in \{L, H\}$ with $\Prob(\theta = H) = \pi$. The worker
chooses education level $e \geq 0$ (the ``signal''). An employer observes $e$ and offers
wage $w(e)$. Payoffs:
\begin{align*}
    u_S(\theta, e, w) &= w - c(\theta, e), \quad \text{with } c(H, e) < c(L, e)
    \text{ for } e > 0, \\
    u_R(\theta, e, w) &= \theta - w,
\end{align*}
where $c(\theta, e)$ is the cost of education (cheaper for high-ability workers:
\emph{single-crossing property}).

\textbf{Separating equilibrium.}
In a separating PBE, $e_H \neq e_L$ and the receiver's beliefs at $e_H$ are
$\mu(H | e_H) = 1$, at $e_L$ are $\mu(H | e_L) = 0$ (by Bayes' rule). The receiver offers
$w(e_H) = H$ and $w(e_L) = L$.

For this to be an equilibrium, each type must prefer its own signal:
\begin{align*}
    H - c(H, e_H) &\geq L - c(H, e_L) \quad \text{(type $H$ prefers $e_H$)}, \\
    L - c(L, e_L) &\geq H - c(L, e_H) \quad \text{(type $L$ prefers $e_L$)}.
\end{align*}
The single-crossing property ensures that such $e_H, e_L$ exist (typically $e_L = 0$ and
$e_H$ above a threshold).

\textbf{Pooling equilibrium.}
Both types choose $e^* = 0$. The receiver's belief at $e^*$ is the prior:
$\mu(H | e^*) = \pi$, so $w(e^*) = \pi H + (1-\pi) L$. Off-equilibrium beliefs at
$e \neq e^*$ can be chosen to deter deviation, e.g., $\mu(H | e) = 0$ for all $e \neq e^*$,
giving $w(e) = L$. Then no type wants to deviate if
$\pi H + (1-\pi)L - c(\theta, 0) \geq L - c(\theta, e)$ for all $e > 0$, which holds
when the pooling wage is sufficiently attractive.
\end{example}

\begin{proposition}[PBE Existence in Finite Signaling Games]
\label{prop:pbe-signaling-existence}
Every finite signaling game (finite $\Theta$, $M$, $A$) has a perfect Bayesian equilibrium.
\end{proposition}

\begin{proof}
A finite signaling game is a finite extensive-form game. The sequential equilibrium
existence theorem of Kreps and Wilson (1982) guarantees existence of a sequential
equilibrium, which is a refinement of PBE. In particular, every sequential equilibrium
is a PBE, so PBE exists.

Alternatively, one can construct a PBE directly: consider all possible sender strategies
$\sigma_S : \Theta \to \Delta(M)$. For each such strategy, compute the receiver's
beliefs via Bayes' rule at information sets reached with positive probability, and assign
arbitrary beliefs at unreached information sets. The receiver's optimal response defines
a best-response map. The sender's optimal message choice, given the receiver's response
function, defines another best-response map. By Kakutani's fixed point theorem (or Nash's
theorem applied to the agent normal form), a fixed point---hence a PBE---exists.
\end{proof}


%% ============================================================
%% SECTION 8: Potential Games
%% ============================================================
\section{Potential Games}
\label{sec:potential-games}

\subsection{Definitions}

Potential games form a special class of games where all players' incentives can be captured
by a single ``potential function.'' This structure has powerful implications for existence
of pure equilibria and convergence of learning dynamics.

\begin{definition}[Exact Potential Game]
\label{def:exact-potential}
A finite game $G = (N, (A_i)_{i \in N}, (u_i)_{i \in N})$ is an \emph{exact potential
game} if there exists a function $\Phi : A \to \R$ (the \emph{potential function}) such
that for every player $i$, every $a_{-i} \in A_{-i}$, and every $a_i, a_i' \in A_i$:
\[
    u_i(a_i, a_{-i}) - u_i(a_i', a_{-i}) = \Phi(a_i, a_{-i}) - \Phi(a_i', a_{-i}).
\]
\end{definition}

\begin{definition}[Ordinal Potential Game]
\label{def:ordinal-potential}
A finite game $G$ is an \emph{ordinal potential game} if there exists $\Phi : A \to \R$
such that for every player $i$, every $a_{-i}$, and every $a_i, a_i'$:
\[
    u_i(a_i, a_{-i}) > u_i(a_i', a_{-i}) \iff \Phi(a_i, a_{-i}) > \Phi(a_i', a_{-i}).
\]
\end{definition}

\begin{remark}
Every exact potential game is an ordinal potential game. The converse is false: ordinal
potentials only preserve the \emph{sign} of payoff differences, not their magnitude.
\end{remark}

\subsection{Existence of Pure Nash Equilibria}

\begin{theorem}[Potential Games Have Pure NE]
\label{thm:potential-pure-ne}
Every finite exact potential game (and every finite ordinal potential game) has at least one
pure Nash equilibrium.
\end{theorem}

\begin{proof}
Let $\Phi : A \to \R$ be the potential function. Since $A$ is finite, $\Phi$ attains its
maximum at some $a^* \in A$:
\[
    a^* \in \argmax_{a \in A} \Phi(a).
\]
We claim $a^*$ is a pure Nash equilibrium.

For any player $i$ and any deviation $a_i' \in A_i$:
\[
    \Phi(a_i^*, a_{-i}^*) \geq \Phi(a_i', a_{-i}^*)
\]
by the maximality of $a^*$. By the potential property:
\[
    u_i(a_i^*, a_{-i}^*) - u_i(a_i', a_{-i}^*)
    = \Phi(a_i^*, a_{-i}^*) - \Phi(a_i', a_{-i}^*) \geq 0.
\]
Thus $a_i^*$ is a best response to $a_{-i}^*$ for every player $i$, so $a^*$ is a NE.

For ordinal potential games, the same argument works: $\Phi(a_i^*, a_{-i}^*) \geq
\Phi(a_i', a_{-i}^*)$ implies $u_i(a_i^*, a_{-i}^*) \geq u_i(a_i', a_{-i}^*)$ by the
ordinal potential property (since if the strict inequality failed, i.e.,
$u_i(a_i', a_{-i}^*) > u_i(a_i^*, a_{-i}^*)$, then
$\Phi(a_i', a_{-i}^*) > \Phi(a_i^*, a_{-i}^*)$, contradicting maximality).
\end{proof}

\subsection{Convergence of Best-Response Dynamics}

\begin{definition}[Best-Response Dynamics]
\label{def:br-dynamics}
In \emph{(asynchronous) best-response dynamics}, players take turns updating their actions.
At each step, one player $i$ is selected (deterministically or randomly) and switches to a
best response to the current actions of all other players:
\[
    a_i^{t+1} \in \argmax_{a_i \in A_i} u_i(a_i, a_{-i}^t),
    \qquad a_j^{t+1} = a_j^t \text{ for } j \neq i.
\]
\end{definition}

\begin{theorem}[Finite Improvement Property]
\label{thm:fip}
In every finite exact potential game, best-response dynamics converge to a pure Nash
equilibrium in a finite number of steps.
\end{theorem}

\begin{proof}
We show that $\Phi$ strictly increases at each step where a player's action changes, hence
the dynamics must terminate (since $A$ is finite and $\Phi$ cannot increase indefinitely).

Suppose at step $t$, player $i$ switches from $a_i^t$ to $a_i^{t+1} \neq a_i^t$ with
$a_i^{t+1} \in \argmax_{a_i} u_i(a_i, a_{-i}^t)$. Since $a_i^{t+1}$ is a strictly
improving move (otherwise the player would not switch), we have:
\[
    u_i(a_i^{t+1}, a_{-i}^t) > u_i(a_i^t, a_{-i}^t).
\]
By the exact potential property:
\[
    \Phi(a_i^{t+1}, a_{-i}^t) - \Phi(a_i^t, a_{-i}^t)
    = u_i(a_i^{t+1}, a_{-i}^t) - u_i(a_i^t, a_{-i}^t) > 0.
\]

Thus $\Phi(a^{t+1}) > \Phi(a^t)$ whenever an action changes. Since $A$ is finite, the
sequence $\Phi(a^1), \Phi(a^2), \ldots$ is strictly increasing and takes values in the
finite set $\{\Phi(a) : a \in A\}$, so it must terminate. Termination occurs when no player
wants to change, i.e., at a pure Nash equilibrium.

\emph{Remark on step count:} The number of steps is bounded by $|A| = \prod_i |A_i|$
(since each action profile can appear at most once in the improvement path), but this
worst-case bound is typically very loose.
\end{proof}

\begin{corollary}
\label{cor:better-response}
The finite improvement property also holds for \emph{better-response dynamics} (where
player $i$ switches to any action that improves her payoff, not necessarily the best
response), since $\Phi$ still strictly increases.
\end{corollary}

\begin{proof}
If $u_i(a_i', a_{-i}^t) > u_i(a_i^t, a_{-i}^t)$, then
$\Phi(a_i', a_{-i}^t) > \Phi(a_i^t, a_{-i}^t)$ by the potential property. The same
termination argument applies.
\end{proof}

\subsection{Congestion Games}

The most important class of potential games in applications is congestion games.

\begin{definition}[Congestion Game]
\label{def:congestion-game}
A \emph{congestion game} is defined by $(N, E, (A_i)_{i \in N}, (d_e)_{e \in E})$ where:
\begin{itemize}
    \item $N = [n]$ is the player set.
    \item $E$ is a finite set of \emph{resources} (or \emph{facilities}).
    \item $A_i \subseteq 2^E$ is a collection of subsets of resources available to player
    $i$. Each action $a_i \in A_i$ is a set of resources that player $i$ uses.
    \item $d_e : \{1, \ldots, n\} \to \R$ is the \emph{delay function} (or cost function)
    for resource $e$, depending on the number of users.
\end{itemize}
The payoff (negative cost) to player $i$ under action profile $a = (a_1, \ldots, a_n)$ is:
\[
    u_i(a) = -\sum_{e \in a_i} d_e\bigl(n_e(a)\bigr),
\]
where $n_e(a) = |\{j \in N : e \in a_j\}|$ is the congestion (number of users) of
resource $e$.
\end{definition}

\begin{theorem}[Congestion Games are Exact Potential Games (Rosenthal)]
\label{thm:rosenthal}
Every congestion game is an exact potential game with potential function
\[
    \Phi(a) = -\sum_{e \in E} \sum_{k=1}^{n_e(a)} d_e(k).
\]
\end{theorem}

\begin{proof}
We must verify that for any player $i$, any $a_{-i}$, and any $a_i, a_i' \in A_i$:
\[
    u_i(a_i, a_{-i}) - u_i(a_i', a_{-i}) = \Phi(a_i, a_{-i}) - \Phi(a_i', a_{-i}).
\]

Consider the left side. Let $n_e = n_e(a_i, a_{-i})$ denote the congestion when player $i$
uses $a_i$, and $n_e' = n_e(a_i', a_{-i})$ when player $i$ uses $a_i'$. We partition the
resources into three groups:
\begin{enumerate}[label=(\roman*)]
    \item $e \in a_i \cap a_i'$: used by $i$ under both actions. Then $n_e = n_e'$, so
    $d_e(n_e)$ cancels in both $u_i$ and $\Phi$ differences.
    \item $e \in a_i \setminus a_i'$: used by $i$ under $a_i$ but not $a_i'$. Then
    $n_e = n_e' + 1$ (one more user under $a_i$). The contribution to $u_i(a_i, a_{-i})
    - u_i(a_i', a_{-i})$ is $-d_e(n_e) = -d_e(n_e' + 1)$. The contribution to
    $\Phi(a_i, a_{-i}) - \Phi(a_i', a_{-i})$ is $-\sum_{k=1}^{n_e'+1} d_e(k)
    + \sum_{k=1}^{n_e'} d_e(k) = -d_e(n_e'+1)$. These are equal.
    \item $e \in a_i' \setminus a_i$: used by $i$ under $a_i'$ but not $a_i$. Then
    $n_e' = n_e + 1$. The contribution to $u_i(a_i, a_{-i}) - u_i(a_i', a_{-i})$ is
    $0 - (-d_e(n_e')) = d_e(n_e + 1)$. The contribution to
    $\Phi(a_i, a_{-i}) - \Phi(a_i', a_{-i})$ is $-\sum_{k=1}^{n_e} d_e(k)
    + \sum_{k=1}^{n_e+1} d_e(k) = d_e(n_e+1)$. These are equal.
\end{enumerate}
Summing over all resources, $u_i(a_i, a_{-i}) - u_i(a_i', a_{-i})
= \Phi(a_i, a_{-i}) - \Phi(a_i', a_{-i})$.
\end{proof}

\begin{corollary}
\label{cor:congestion-ne}
Every finite congestion game has a pure Nash equilibrium, and best-response dynamics
converge to one in finitely many steps.
\end{corollary}

\begin{proof}
Immediate from Theorems~\ref{thm:rosenthal}, \ref{thm:potential-pure-ne},
and~\ref{thm:fip}.
\end{proof}

\begin{example}[Network Routing Game]
\label{ex:routing}
Consider a network with source $s$ and destination $d$, and $n$ users each choosing a path
from $s$ to $d$. Each edge $e$ has delay $d_e(k)$ depending on congestion $k$. This is a
congestion game: resources are edges, each player's action set is the set of $s$-$d$ paths,
and the total delay on a path is $\sum_{e \in \text{path}} d_e(n_e)$. By
Corollary~\ref{cor:congestion-ne}, a pure NE exists and can be found by best-response
dynamics.
\end{example}

\begin{example}[Resource Allocation]
\label{ex:resource-allocation}
Consider $n$ agents each choosing one of $m$ machines to process a job. Machine $e$ has
processing time $d_e(k) = k \cdot \ell_e$ (linear congestion) where $\ell_e$ is the base
load. Each agent's cost is the processing time on her chosen machine. This is a congestion
game with $A_i = E = \{1, \ldots, m\}$ (each action uses a single resource). The potential
function is $\Phi(a) = -\sum_{e=1}^m \ell_e \cdot \frac{n_e(n_e+1)}{2}$, and pure NE exist
with convergent best-response dynamics.
\end{example}

\begin{rlconnection}[Potential Games and Independent Learning in MARL]
Potential games are among the few game classes where independent learners (agents running
separate RL algorithms without coordination) are guaranteed to converge:
\begin{itemize}
    \item \textbf{Best-response dynamics} correspond to each agent independently computing
    optimal policies given fixed opponent policies. The finite improvement property ensures
    convergence.
    \item \textbf{Log-linear learning:} If each agent plays action $a_i$ with probability
    proportional to $\exp(\beta \cdot u_i(a_i, a_{-i}))$ (Boltzmann exploration), the
    joint process is a Markov chain that converges to the potential maximizer as
    $\beta \to \infty$.
    \item \textbf{Multi-agent reward shaping:} If a team reward function $\Phi$ can serve
    as a potential for individual reward functions $u_i$, then shaping individual rewards
    via $\Phi$ makes the game a potential game, guaranteeing convergence.
\end{itemize}
This is the theoretical basis for \emph{difference rewards} and \emph{Wonderful Life
Utility} in cooperative MARL.
\end{rlconnection}


%% ============================================================
%% SECTION 9: Connections to Multi-Agent RL
%% ============================================================
\section{Connections to Multi-Agent RL}
\label{sec:marl-connections}

We consolidate the connections between the game-theoretic results of this lesson and
modern multi-agent reinforcement learning.

\begin{rlconnection}[Independent Learners in Repeated Games]
When $n$ agents independently run Q-learning or policy gradient in a repeated game, each
agent treats the others as part of the environment. The folk theorems
(Theorems~\ref{thm:nash-folk} and~\ref{thm:spe-folk}) imply that the set of stable
outcomes is vast---any feasible, individually rational payoff can be a fixed point. This
explains the well-known instability and sensitivity to hyperparameters of independent
learners in general-sum games. Convergence guarantees for independent Q-learning exist
only in restrictive settings: two-player zero-sum games (minimax-Q), potential games, and
certain classes of team games.
\end{rlconnection}

\begin{rlconnection}[Self-Play and Fictitious Play]
Self-play---where an agent trains against copies of itself---is a form of fictitious play.
In zero-sum games, Robinson's theorem (Theorem~\ref{thm:robinson}) guarantees that the
empirical strategy profile converges to NE. Modern algorithms extend this:
\begin{itemize}
    \item \textbf{Neural Fictitious Self-Play (NFSP):} Uses neural networks to approximate
    best responses and average strategies, converging to approximate NE in large zero-sum
    games.
    \item \textbf{PSRO (Policy-Space Response Oracles):} Maintains a population of
    policies and computes a best response to the meta-mixture, directly implementing
    fictitious play in policy space.
    \item \textbf{Double Oracle:} Iteratively adds best responses to a restricted game,
    converging to the NE of the full game.
\end{itemize}
The regret bound of MWU (Theorem~\ref{thm:mwu-regret}) provides the convergence rate:
$O(1/\sqrt{T})$ for last-iterate convergence in zero-sum games with optimistic updates.
\end{rlconnection}

\begin{rlconnection}[Opponent Modeling as Bayesian Inference]
The Bayesian game framework (Section~\ref{sec:bayesian-games}) provides the foundation for
opponent modeling in MARL:
\begin{enumerate}[label=(\roman*)]
    \item \textbf{Type inference:} An agent maintains a belief $p(\theta_{-i} | h^t)$ over
    opponent types, updated via Bayes' rule as actions are observed.
    \item \textbf{Bayesian policy optimization:} The agent's policy maximizes expected
    return under the current belief: $\max_{\pi_i} \E_{\theta_{-i} \sim p(\cdot | h^t)}
    [R_i(\pi_i, \pi_{-i}(\theta_{-i}))]$.
    \item \textbf{Exploration-exploitation:} The agent may take actions to learn about
    opponents' types (information-gathering exploration), connecting to the theory of
    Bayesian exploration in MDPs.
\end{enumerate}
\end{rlconnection}

\begin{rlconnection}[Potential-Based Reward Shaping in Multi-Agent Settings]
Ng, Harada, and Russell's (1999) potential-based reward shaping theorem states that
adding a shaping reward $F(s, s') = \gamma \phi(s') - \phi(s)$ preserves the optimal
policy in single-agent MDPs. In multi-agent settings, this connects to potential games:
\begin{itemize}
    \item If the joint reward function serves as a potential for individual rewards,
    the multi-agent system is a potential game and independent learning converges.
    \item The \emph{difference reward} $D_i(a) = R(a) - R(a_{-i}, c_i)$ (where $c_i$ is a
    default action for agent $i$) is designed so that $D_i(a_i, a_{-i}) - D_i(a_i', a_{-i})
    = R(a_i, a_{-i}) - R(a_i', a_{-i})$, making $R$ an exact potential for the game with
    payoffs $D_i$.
    \item This theoretical property ensures that maximizing individual difference rewards
    aligns with maximizing the global objective, providing a principled approach to
    cooperative MARL.
\end{itemize}
\end{rlconnection}


%% ============================================================
%% SECTION 10: Exercises
%% ============================================================
\section{Exercises}
\label{sec:exercises}

\begin{exercise}[Finitely Repeated Prisoner's Dilemma]
\label{ex:finite-pd}
Consider the Prisoner's Dilemma from Example~\ref{ex:repeated-pd}, repeated $T = 100$
times. Show that the unique SPE is mutual defection in every period, regardless of $T$.
Explain why this result is ``paradoxical'' and how it motivates infinitely repeated games.
\end{exercise}

\begin{exercise}[Grim Trigger in a Different Game]
\label{ex:grim-trigger-general}
Consider the stage game:
\[
\begin{array}{c|ccc}
  & L & C & R \\ \hline
T & 6,6 & 0,8 & 0,0 \\
M & 8,0 & 4,4 & 0,0 \\
B & 0,0 & 0,0 & 2,2
\end{array}
\]
\begin{enumerate}[label=(\alph*)]
    \item Find all pure Nash equilibria of the stage game.
    \item Compute the minimax values $\underline{v}_1$ and $\underline{v}_2$.
    \item Show that the payoff $(6,6)$ (from $(T,L)$) can be sustained as an SPE of the
    infinitely repeated game using grim trigger with punishment $(B,R)$. Find the minimum
    discount factor.
    \item Can the payoff $(7,7)$ be sustained? If so, construct a strategy profile.
\end{enumerate}
\end{exercise}

\begin{exercise}[Tit-for-Tat Is Not SPE]
\label{ex:tft-not-spe}
Show that in the infinitely repeated Prisoner's Dilemma (Example~\ref{ex:repeated-pd}),
the tit-for-tat strategy profile is a Nash equilibrium for $\delta \geq 1/3$ but is
\emph{not} a subgame perfect equilibrium. (Hint: Consider the subgame after $(D, C)$.)
\end{exercise}

\begin{exercise}[Computing BNE]
\label{ex:compute-bne}
Consider a two-player Bayesian game where player 1 has type $\theta_1 \in \{H, L\}$ with
$\Prob(\theta_1 = H) = 1/2$, and player 2 has no private information. Actions:
$A_1 = A_2 = \{U, D\}$. Payoffs (for each type of player 1):
\[
\text{Type } H: \quad
\begin{array}{c|cc}
  & L & R \\ \hline
U & 2,1 & 0,0 \\
D & 0,0 & 1,2
\end{array}
\qquad \text{Type } L: \quad
\begin{array}{c|cc}
  & L & R \\ \hline
U & 0,0 & 1,2 \\
D & 2,1 & 0,0
\end{array}
\]
Find all Bayesian Nash equilibria (pure and mixed).
\end{exercise}

\begin{exercise}[Second-Price Auction as Bayesian Game]
\label{ex:second-price}
In a second-price sealed-bid auction with $n$ bidders having independent private values
$v_i \sim F$ (where $F$ is a continuous CDF on $[0, \bar{v}]$), show that bidding
$b_i = v_i$ (truthful bidding) is a dominant strategy, hence a BNE. Compare with the
first-price auction BNE from Example~\ref{ex:first-price-auction}.
\end{exercise}

\begin{exercise}[Verifying a Potential Function]
\label{ex:verify-potential}
Consider the two-player game with $A_1 = A_2 = \{X, Y\}$ and payoffs:
\[
\begin{array}{c|cc}
  & X & Y \\ \hline
X & 3,2 & 1,1 \\
Y & 2,1 & 4,3
\end{array}
\]
\begin{enumerate}[label=(\alph*)]
    \item Show this is an exact potential game by constructing the potential function $\Phi$.
    \item Find all pure Nash equilibria using $\Phi$.
    \item Verify that best-response dynamics converge from any initial profile.
\end{enumerate}
\end{exercise}

\begin{exercise}[MWU Convergence Rate]
\label{ex:mwu-rate}
Consider the zero-sum game Rock-Paper-Scissors:
\[
\begin{array}{c|ccc}
  & R & P & S \\ \hline
R & 0 & -1 & 1 \\
P & 1 & 0 & -1 \\
S & -1 & 1 & 0
\end{array}
\]
\begin{enumerate}[label=(\alph*)]
    \item What is the unique NE and the value of the game?
    \item If both players use MWU with $\eta = \sqrt{\ln 3 / T}$, bound the duality gap
    $\max_{a_1} g(a_1, \bar{\alpha}_2^T) - \min_{a_2} g(\bar{\alpha}_1^T, a_2)$ after
    $T = 1000$ rounds.
    \item How many rounds are needed for the duality gap to be below $0.01$?
\end{enumerate}
\end{exercise}

\begin{exercise}[Folk Theorem Application]
\label{ex:folk-application}
Consider the two-player stage game:
\[
\begin{array}{c|cc}
  & C & D \\ \hline
C & 5,5 & 0,7 \\
D & 7,0 & 2,2
\end{array}
\]
\begin{enumerate}[label=(\alph*)]
    \item Compute the feasible payoff set $V$ and the minimax values.
    \item Apply the Nash folk theorem to determine which payoff vectors are achievable as
    NE payoffs for large $\delta$.
    \item For the payoff $(5,5)$, find the minimum discount factor under grim trigger.
    \item For the payoff $(4,4)$, describe a strategy profile achieving it as an SPE and
    find the minimum $\delta$.
\end{enumerate}
\end{exercise}

\begin{exercise}[Congestion Game with Three Players]
\label{ex:congestion-three}
Three players share two resources $\{A, B\}$. Each player chooses exactly one resource.
Delay functions: $d_A(k) = k$ and $d_B(k) = 2k - 1$.
\begin{enumerate}[label=(\alph*)]
    \item Write the Rosenthal potential function $\Phi$.
    \item Enumerate all action profiles and compute $\Phi$ for each.
    \item Find all pure NE by maximizing $\Phi$ (equivalently, minimizing $-\Phi$).
    \item Verify that your NE are indeed equilibria by checking no player wants to deviate.
\end{enumerate}
\end{exercise}

\begin{exercise}[Separating vs.\ Pooling Equilibria]
\label{ex:signaling}
Consider a signaling game with two sender types $\theta \in \{G, B\}$ (good, bad),
$\Prob(\theta = G) = 1/3$, two messages $m \in \{H, L\}$ (high, low effort), and two
receiver actions $a \in \{A, R\}$ (accept, reject). Payoffs:
\begin{align*}
    u_S(G, H, A) &= 3, \quad u_S(G, H, R) = -1, \quad u_S(G, L, A) = 4, \quad
    u_S(G, L, R) = 0, \\
    u_S(B, H, A) &= 3, \quad u_S(B, H, R) = -2, \quad u_S(B, L, A) = 4, \quad
    u_S(B, L, R) = 1, \\
    u_R(\theta, m, A) &= \begin{cases} 1 & \text{if } \theta = G, \\
    -1 & \text{if } \theta = B, \end{cases} \quad
    u_R(\theta, m, R) = 0.
\end{align*}
\begin{enumerate}[label=(\alph*)]
    \item Does a separating PBE exist where type $G$ sends $H$ and type $B$ sends $L$?
    \item Does a pooling PBE exist where both types send $L$?
    \item Does a pooling PBE exist where both types send $H$?
\end{enumerate}
\end{exercise}


%% ============================================================
%% SECTION 11: Summary of Key Results
%% ============================================================
\section{Summary of Key Results}
\label{sec:summary}

\begin{center}
\renewcommand{\arraystretch}{1.4}
\begin{tabular}{@{}p{4.2cm}p{8cm}l@{}}
\toprule
\textbf{Result} & \textbf{Statement} & \textbf{Ref.} \\
\midrule

Unraveling Theorem &
Unique stage-game NE $\Rightarrow$ unique SPE of $G^T$: play NE every period &
Thm~\ref{thm:unraveling} \\

Nash Folk Theorem &
Any feasible, strictly individually rational payoff is a NE payoff for $\delta$ near 1 &
Thm~\ref{thm:nash-folk} \\

SPE Folk Theorem &
Same as Nash folk theorem but for SPE (under full dimensionality condition) &
Thm~\ref{thm:spe-folk} \\

Fictitious Play ($2 \times 2$) &
Empirical frequencies converge to NE in $2 \times 2$ zero-sum games &
Thm~\ref{thm:fp-2x2} \\

Robinson's Theorem &
Fictitious play converges to NE in all finite two-player zero-sum games &
Thm~\ref{thm:robinson} \\

MWU Regret Bound &
External regret $\leq 2\sqrt{T \ln |A_i|}$; hence MWU is no-regret &
Thm~\ref{thm:mwu-regret} \\

No-Regret $\Rightarrow$ NE &
If both players use no-regret, time-avg.\ play converges to NE in zero-sum &
Thm~\ref{thm:no-regret-ne} \\

BNE Existence &
Every finite Bayesian game has a BNE in mixed strategies &
Thm~\ref{thm:bne-existence} \\

PBE in Signaling Games &
Every finite signaling game has a PBE &
Prop~\ref{prop:pbe-signaling-existence} \\

Potential $\Rightarrow$ Pure NE &
Every finite potential game has a pure Nash equilibrium &
Thm~\ref{thm:potential-pure-ne} \\

Finite Improvement &
Best-response dynamics converge in finite steps in potential games &
Thm~\ref{thm:fip} \\

Rosenthal's Theorem &
Every congestion game is an exact potential game &
Thm~\ref{thm:rosenthal} \\

\bottomrule
\end{tabular}
\end{center}

\medskip

\begin{center}
\renewcommand{\arraystretch}{1.4}
\begin{tabular}{@{}p{3.5cm}p{8.5cm}@{}}
\toprule
\textbf{Game Class} & \textbf{MARL Implication} \\
\midrule

Repeated games (general-sum) &
Folk theorems $\Rightarrow$ many equilibria; independent learners may converge to any of them; sensitive to initialization and hyperparameters \\

Zero-sum games &
Minimax + no-regret $\Rightarrow$ convergent self-play (NFSP, PSRO, CFR) \\

Bayesian games &
Opponent modeling as Bayesian inference over types; principled exploration \\

Potential games &
Independent best-response / Q-learning converges; difference rewards make cooperative MARL a potential game \\

Congestion games &
Routing, resource allocation: pure NE exist, local search works \\

\bottomrule
\end{tabular}
\end{center}

\end{document}
